\chapter{Portfolio Manager and Pod Analyst}\label{ch:in:portfolio-manager-and-pod-analyst}

One of the largest multi-manager platforms describes its \omterm{def:in:portfolio-manager-and-pod-analyst:roles}{portfolio managers}' pay in its filing with the US regulator:
a percentage of the profits the manager earned in the preceding year, measured without regard to other managers, with
losses carried forward, a salary that is an advance against the payout, and no clawback. The platform also allocates and
reallocates capital among its managers, and can take it away. The deal is a call option on the manager's own skill with
a knock-out: the manager keeps a share of the upside, bears no share of a loss beyond a lost payout, and loses the job
if the loss goes too deep. Its value depends on the Sharpe ratio the manager actually has, which neither side knows at
the start. This chapter describes the \omterm{def:in:portfolio-manager-and-pod-analyst:roles}{portfolio manager}'s and the analyst's jobs, values the deal, and asks what a track
record can prove.

\begin{center}\footnotesize
\begin{tabular}{@{}p{2.4cm}p{3.3cm}p{3.3cm}p{3.3cm}@{}}
\toprule
\multicolumn{4}{@{}l}{\textbf{Role cards: the pod}}\\
\midrule
 & \omterm{def:in:portfolio-manager-and-pod-analyst:roles}{portfolio manager} & \omterm{def:in:portfolio-manager-and-pod-analyst:roles}{sub-portfolio manager} & \omterm{def:in:portfolio-manager-and-pod-analyst:roles}{pod analyst}\\
\midrule
owns & the pod's P\&L & a slice of it & none; supports the manager's\\
paid by & a share of the pod's profit & a share agreed with the manager & salary and a discretionary bonus\\
reports to & the platform's head of strategy & the \omterm{def:in:portfolio-manager-and-pod-analyst:roles}{portfolio manager} & the \omterm{def:in:portfolio-manager-and-pod-analyst:roles}{portfolio manager}\\
codes in filings & 11-3031, 13-2099.01 & 13-2051, 13-2099.01 & 13-2051\\
taught in & Book~8, ch.~28; Book~16, ch.~3 & Book~16, ch.~3 & Book~7, ch.~1\\
\bottomrule
\end{tabular}
\end{center}

\begin{definition}[Portfolio manager, sub-portfolio manager, pod analyst]\label{def:in:portfolio-manager-and-pod-analyst:roles}
A \emph{portfolio manager}\index{portfolio manager} is the person who decides the positions of a portfolio within the
limits of an owner of the capital and answers for its result. On a platform, a \emph{sub-portfolio
manager}\index{sub-portfolio manager} runs a slice of a portfolio manager's capital with its own limits and a share of
its own profit, inside the manager's team. A \emph{pod analyst}\index{pod analyst} works for a portfolio manager on
research, data and position monitoring, without capital of their own.
\end{definition}

\section{The deal: capital, payout, limits and costs}

The platform's brochure is the best public description of the deal's structure; its numbers are not published. The
structure has four parts.
\begin{itemize}
\item \textbf{Capital.} The platform allocates capital to the pod (Book~8, chapter~28) and ``allocating and
reallocating'' it is its main tool: capital grows after good results and shrinks, or goes, after bad ones.
\item \textbf{Payout.} The manager is paid a share of the pod's own profit, at the payout rate of Book~16, chapter~3,
``without taking into account the performance of other Portfolio Managers''; ``the losses are carried forward and past
losses must be made up before performance-based compensation becomes payable''; the salary is ``generally treated as an
advance''; and managers ``with positive performance will receive performance-based compensation even if'' the funds lose.
\item \textbf{Limits.} Risk limits and drawdown limits (Book~8, chapter~28); the de-risking ladder of Book~16, chapter~8,
cuts capital at a first drawdown and stops the pod at a second.
\item \textbf{Costs.} The pod's own costs (data, staff, the manager's team's pay) are charged against its profit before
the payout, and the investors bear them through the pass-through fee of Book~16, chapter~3.
\end{itemize}
The brochure also records that the platform has agreed to ``guarantee'' a level of pay for a year or years, and to
replace pay that a manager forfeited on leaving a previous employer: the buyouts of chapter~13.

\begin{dated}{2026-08}{A platform's deal with its portfolio managers}\label{dat:in:portfolio-manager-and-pod-analyst:deal}
Millennium Management's Form ADV Part 2A of 25 August 2026: pay ``generally determined as a percentage of profits earned
by such Portfolio Manager during the preceding calendar year''; losses carried forward; salary an advance against the
payout; ``generally no \textquotedbl carryback\textquotedbl{} or \textquotedbl clawback\textquotedbl''; payout due even
when the funds lose; guarantees and buyouts agreed at times; some managers' teams include members ``who may also be
\textquotedbl \omterm{def:in:portfolio-manager-and-pod-analyst:roles}{portfolio managers}\textquotedbl''. No payout rate or drawdown threshold is published.
\end{dated}

\section{The deal as an option with a knock-out}

\begin{method}[The portfolio manager's deal]\label{met:in:portfolio-manager-and-pod-analyst:deal}
Let the pod earn daily returns with annual Sharpe ratio $S$ and volatility $\sigma$ on capital $K$. Each year-end the
manager is paid $\max\bigl(w,\ \rho\max(P_t-L_{t-1},0)K\bigr)$, where $P_t$ is the year's profit, $L_{t-1}$ the loss
carried forward, $\rho$ the payout rate and $w$ the salary; $L_t=\max(L_{t-1}-P_t,0)$. A drawdown from the tenure's
peak beyond $c$ halves the capital for the rest of the tenure; beyond $s$ it ends the job, with the salary paid to that
day and nothing after. The payoff is convex in the year's profit (a call) and truncated by the stop (a knock-out).
\end{method}

\begin{example}[A three-year tenure]\label{ex:in:portfolio-manager-and-pod-analyst:deal}
Take illustrative terms: \$500 million of capital, volatility 6\% a year (\$30 million), payout 15\%, salary \$0.5
million, capital halved at a 5\% drawdown and the pod stopped at 7.5\%, over three years, against a salaried
alternative of \$1.5 million a year. With no skill ($S=0$) the manager's expected three-year pay is \$3.18 million
(median \$1.50 million) and the pod is stopped in 57.3\% of tenures. At $S=1$ the expected pay is \$10.48 million and the
stop probability 15.2\%; at $S=2$, \$23.36 million and 1.9\%. The expected pay reaches the alternative's \$4.5 million
at $S\approx0.26$, where the pod is stopped in about 43\% of tenures (\cref{fig:in:portfolio-manager-and-pod-analyst:deal,fig:in:portfolio-manager-and-pod-analyst:stop}).
\end{example}

\begin{omfigure}
\begin{tikzpicture}
\begin{axis}[width=10.5cm,height=5.6cm,xmin=0,xmax=2,ymin=0,ymax=40,xlabel={\small true annual Sharpe ratio of the pod},
  ylabel={\small three-year pay, \$ million},tick label style={font=\footnotesize},grid=major,grid style={gray!20},
  table/col sep=comma,legend style={font=\scriptsize,at={(0.5,-0.3)},anchor=north,legend cell align=left,
  legend columns=2,/tikz/every even column/.append style={column sep=8pt}}]
\addplot[name path=hi,draw=none,forget plot] table[x=sr,y=p90]{figdata/industry/22-portfolio-manager-and-pod-analyst/deal.csv};
\addplot[name path=lo,draw=none,forget plot] table[x=sr,y=p10]{figdata/industry/22-portfolio-manager-and-pod-analyst/deal.csv};
\addplot[omDef!20] fill between[of=hi and lo];
\addplot[omDef,thick,mark=*] table[x=sr,y=mean]{figdata/industry/22-portfolio-manager-and-pod-analyst/deal.csv};
\addplot[omThm,thick,dashed,mark=square*,mark options={solid}] table[x=sr,y=ce]{figdata/industry/22-portfolio-manager-and-pod-analyst/deal.csv};
\addplot[black!60,thick,domain=0:2] {4.5};
\legend{10th to 90th percentile,expected pay,certainty equivalent (RRA 3),salaried alternative}
\end{axis}
\end{tikzpicture}

\omcaption{The \omterm{def:in:portfolio-manager-and-pod-analyst:roles}{portfolio manager}'s three-year pay against the pod's true Sharpe ratio, with the example's illustrative
terms: mean, 10th--90th percentile band, and the certainty equivalent for relative risk aversion 3 and \$5 million of
other wealth, against a \$1.5 million-a-year salaried alternative. 20\,000 tenures a point. Data:
\texttt{in\_pm.curve}.}\label{fig:in:portfolio-manager-and-pod-analyst:deal}
\end{omfigure}

The deal pays for luck. With no skill at all, its expected value is \$3.18 million over three years, about seven-tenths
of the alternative, because the manager keeps a share of good years and gives back nothing in bad ones except the payouts
that carried losses cancel. The investors pay for that option, which is the netting cost of Book~16, chapter~3, and the platform
limits it with the stop: at $S=0$ more than half the tenures end early. For a risk-averse manager the calculation moves:
with relative risk aversion 3 and \$5 million of other wealth, the certainty equivalent reaches the alternative only at
$S\approx0.75$, where the stop probability is about 23\%. The spread of outcomes is wide at every skill
(\cref{fig:in:portfolio-manager-and-pod-analyst:hist}).

\begin{omfigure}
\begin{tikzpicture}
\begin{axis}[width=10.5cm,height=4.8cm,xmin=0,xmax=2,ymin=0,ymax=60,xlabel={\small true annual Sharpe ratio of the pod},
  ylabel={\small tenures stopped, \%},tick label style={font=\footnotesize},grid=major,grid style={gray!20},
  table/col sep=comma]
\addplot[omThm,thick,mark=*] table[x=sr,y=stopped]{figdata/industry/22-portfolio-manager-and-pod-analyst/deal.csv};
\end{axis}
\end{tikzpicture}

\omcaption{Share of three-year tenures that end at the stop, against the pod's true Sharpe ratio, with the example's
ladder (capital halved at a 5\% drawdown, pod stopped at 7.5\%). Data: \texttt{in\_pm.curve}.}\label{fig:in:portfolio-manager-and-pod-analyst:stop}
\end{omfigure}

\begin{omfigure}
\begin{tikzpicture}
\begin{axis}[ybar interval,area legend,width=10.5cm,height=5cm,xmin=0,xmax=40,ymin=0,ymax=60,x tick label as interval=false,
  xtick={0,4,8,12,16,20,24,28,32,36,40},
  xlabel={\small three-year pay, \$ million (bins of 2)},ylabel={\small share of tenures, \%},
  tick label style={font=\footnotesize},grid=major,grid style={gray!20},table/col sep=comma,
  legend style={font=\scriptsize,at={(0.97,0.97)},anchor=north east,legend cell align=left}]
\addplot[fill=black!25,draw=black!50] table[x=left,y=sr0]{figdata/industry/22-portfolio-manager-and-pod-analyst/hist.csv};
\addplot[fill=omDef!40,draw=omDef,fill opacity=0.6] table[x=left,y=sr1]{figdata/industry/22-portfolio-manager-and-pod-analyst/hist.csv};
\legend{no skill ($S=0$),$S=1$}
\end{axis}
\end{tikzpicture}

\omcaption{Distribution of three-year pay at two true Sharpe ratios, with the example's terms (pay above \$40 million put
in the last bin). Without skill most tenures pay little more than salary; with a Sharpe ratio of one the pay is spread
from salary to over \$30 million. Data: \texttt{in\_pm.run}.}\label{fig:in:portfolio-manager-and-pod-analyst:hist}
\end{omfigure}

\section{The analyst: from research to running capital}

A pod is a small team: a \omterm{def:in:portfolio-manager-and-pod-analyst:roles}{portfolio manager}, one or more analysts, sometimes a \omterm{def:in:quant-researcher:def}{quantitative researcher} or developer, and on
larger pods \omterm{def:in:portfolio-manager-and-pod-analyst:roles}{sub-portfolio managers}. The brochure notes that oversight of part of a manager's assets may be given to
members of the team ``who may also be \textquotedbl \omterm{def:in:portfolio-manager-and-pod-analyst:roles}{portfolio managers}\textquotedbl'', and that some managers mainly
``oversee and manage other investment personnel''. The analyst's work is the manager's research: models of the companies
or markets the pod trades, data, the monitoring of positions and risk, and ideas.

The analyst's pay is set by the manager, out of the pod's economics, and is less formulaic than the manager's. The
route up is to run capital: first a slice of the manager's book as a \omterm{def:in:portfolio-manager-and-pod-analyst:roles}{sub-portfolio manager}, with its own record, then a
pod of one's own on the same platform or another. That record is what the next section is about.

\section{Track records and moving}

A manager moves with a track record, and the platform hiring them must judge it. The portability of a track record
(Book~16, chapter~25) is limited by what it can prove. Chapter~17's arithmetic applies: a true Sharpe ratio of 1 needs
3.85 years of daily results to be distinguishable from zero at 95\%, and a three-year record of Sharpe ratio 1 has a
t-statistic of about 1.73. A manager stopped in year two has a record too short to prove or disprove anything, and a
manager who survived three years may have been lucky: at $S=0$, 43\% of the example's tenures survive three years.

The hiring platform therefore reads more than the number: the strategy's capacity and crowding, the risk taken for the
return, how much of the result came from a few trades, whether the process can be explained, and whether the team moves
with the manager. The restrictive covenants of chapter~28 decide when the manager can start, and the buyout of chapter~13
what the move costs the new employer.

\section{Portfolio managers outside the platforms}

The \omterm{def:in:portfolio-manager-and-pod-analyst:roles}{portfolio manager}'s title is older than the platforms and covers different deals.
\begin{itemize}
\item \textbf{At a single-manager hedge fund} (chapter~4) the \omterm{def:in:portfolio-manager-and-pod-analyst:roles}{portfolio manager} may be the founder, paid through the
firm's fees and ownership rather than a payout rate (Book~16, chapter~25, on launching a fund).
\item \textbf{At a systematic fund} (chapter~4) portfolios are run by models; ``\omterm{def:in:portfolio-manager-and-pod-analyst:roles}{portfolio manager}'' may mean the person
who oversees a set of strategies and their risk.
\item \textbf{At an asset manager or pension fund} (chapter~8) the \omterm{def:in:portfolio-manager-and-pod-analyst:roles}{portfolio manager} runs a mandate against a benchmark
and is judged against it; nothing in the mandate works like a platform's stop.
\end{itemize}
The survey's financial managers, whose detailed occupations include investment fund managers (O*NET lists ``Portfolio
Manager'' among their reported titles), earned in the securities industry a median wage of
\$223\,860 in May 2025, with a 90th percentile of \$383\,590; the analysts \$124\,370 and \$250\,000. The survey records
wages only, which for a platform manager are a small part of pay.

\begin{dated}{2025-05}{What portfolio managers and analysts are paid: the public evidence}\label{dat:in:portfolio-manager-and-pod-analyst:pay}
Occupational survey, May 2025, securities industry: financial managers (11-3031) 70\,400 employed, 10th--90th percentile
\$128\,260--383\,590, median \$223\,860; financial and investment analysts (13-2051) 89\,390, median \$124\,370, 90th
percentile \$250\,000. Other investment pools and funds: financial managers median \$215\,690, analysts \$125\,710. The
labour condition applications for portfolio-manager titles at the sourced employers are too few to publish (14 in fiscal
2025). Wages only.
\end{dated}

\begin{omfigure}
\begin{tikzpicture}
\begin{axis}[width=10.5cm,height=4.8cm,xmin=40,xmax=400,ymin=0.4,ymax=4.6,ytick=data,
  yticklabels from table={figdata/industry/22-portfolio-manager-and-pod-analyst/survey.csv}{label},
  yticklabel style={font=\scriptsize},xticklabel style={font=\footnotesize},xlabel={\small annual wage, \$ thousand},
  grid=major,grid style={gray!20},table/col sep=comma]
\addplot[draw=none,mark=none,error bars/.cd,x dir=both,x explicit,error mark=none,error bar style={omDef,line width=1pt}]
  table[x=p50,y=pos,x error plus expr=\thisrow{p90}-\thisrow{p50},x error minus expr=\thisrow{p50}-\thisrow{p10}]
  {figdata/industry/22-portfolio-manager-and-pod-analyst/survey.csv};
\addplot[draw=none,mark=none,error bars/.cd,x dir=both,x explicit,error mark=none,error bar style={omDef!40,line width=7pt}]
  table[x=p50,y=pos,x error plus expr=\thisrow{p75}-\thisrow{p50},x error minus expr=\thisrow{p50}-\thisrow{p25}]
  {figdata/industry/22-portfolio-manager-and-pod-analyst/survey.csv};
\addplot[only marks,mark=|,mark size=5pt,line width=1.5pt,omThm] table[x=p50,y=pos]{figdata/industry/22-portfolio-manager-and-pod-analyst/survey.csv};
\end{axis}
\end{tikzpicture}

\omcaption{Wages of financial managers and of financial and investment analysts in the securities industry and in other
investment pools and funds, May 2025 occupational survey: 10th to 90th percentile (thin), interquartile range (thick),
median (mark). Wages exclude bonuses and payouts. Data: \texttt{data/industry/oews\_roles.csv}, through
\texttt{in\_pm.survey}.}\label{fig:in:portfolio-manager-and-pod-analyst:survey}
\end{omfigure}

\section{Tutorial: valuing the deal}\label{tut:in:portfolio-manager-and-pod-analyst:tut}

\begin{tutorial}
\textbf{Goal.} Value a platform \omterm{def:in:portfolio-manager-and-pod-analyst:roles}{portfolio manager}'s deal as a function of skill, and find the skill at which it beats a
salaried job.
\textbf{End state:} \cref{fig:in:portfolio-manager-and-pod-analyst:deal,fig:in:portfolio-manager-and-pod-analyst:stop,fig:in:portfolio-manager-and-pod-analyst:hist}.
\begin{tutsteps}
\item \textbf{Pods.} Book~16's \texttt{firm.podshop.pods} draws daily returns for independent pods of a given Sharpe
ratio and volatility.
\item \textbf{The deal.} \texttt{firm.roles.pm\_deal} applies the payout, the carried losses, the salary advance and the
ladder over the tenure (\cref{lst:in:portfolio-manager-and-pod-analyst:deal}).
\omcode{code/firm/roles/firm_roles.py}{347}{379}{The portfolio manager's deal over a tenure, with a two-step drawdown
ladder.}\label{lst:in:portfolio-manager-and-pod-analyst:deal}
\item \textbf{Utility.} Chapter~13's \texttt{firm.payoffer.certainty\_equivalent} converts the pay distribution into a
sure amount.
\item \textbf{The crossing.} \texttt{in\_pm.crossing} interpolates the Sharpe ratio at which expected pay, or the
certainty equivalent, reaches the alternative.
\end{tutsteps}
\end{tutorial}

With the example's terms: expected pay reaches \$4.5 million at $S\approx0.26$ (stop probability about 43\%); the
certainty equivalent at $S\approx0.75$ (about 23\%).

\textbf{What to change next.} Let the platform raise capital after good years; add a guarantee in year one; make the
Sharpe ratio uncertain and let the manager learn it from results; charge the pod's costs.

\section{Build: pod cards and the manager's deal}\label{bld:in:portfolio-manager-and-pod-analyst:roles}

\begin{build}
\bfield{Purpose} Add the pod's roles to the registry and model the \omterm{def:in:portfolio-manager-and-pod-analyst:roles}{portfolio manager}'s deal.
\bfield{Interface} \texttt{firm.roles}: the cards \texttt{portfolio manager}, \texttt{sub-portfolio manager},
\texttt{pod analyst}; \texttt{pm\_deal(sr, years, vol, capital, rate, salary, cut, stop, n, rng, days) ->
dict(pay, stopped, stop\_day, pnl)}; on \texttt{firm.podshop.pods}.
\bfield{Rules} Payout on the pod's own profit above carried losses; the salary is an advance; no clawback; the ladder
works on the drawdown from the tenure's peak; terms are the caller's, labelled illustrative.
\bfield{Acceptance tests} \texttt{code/firm/roles/tests/}: a pod that cannot lose pays exactly the rate times its profit;
a stop on the first day pays one day's salary; more skill means fewer stops and more pay.
\bfield{Stretch} Capital that grows with results; a manager who stops trading after a payout-maximising year; the
platform's side: expected payout against the investors' net return.
\end{build}

\begin{omsources}
\item Millennium Management LLC, Form ADV Part 2A brochure, 25 August 2026 (SEC IAPD).
\item Bureau of Labor Statistics, occupational survey, May 2025; chapter~14's LCA tables.
\item Book~16, chapters~3, 8 and~25; Book~8, chapter~28.
\end{omsources}

\section{Exercises}

\begin{exercise}[$\star$]\label{exo:in:portfolio-manager-and-pod-analyst:1}
A pod made \$40 million last year after a \$10 million loss the year before. At a 15\% payout rate and a \$0.5 million
salary, what is the manager paid for last year?
\end{exercise}

\begin{exercise}[$\star$]\label{exo:in:portfolio-manager-and-pod-analyst:2}
From the example, how much does the expected three-year pay rise between $S=0$ and $S=1$, and how much does the stop
probability fall?
\end{exercise}

\begin{exercise}[$\star$]\label{exo:in:portfolio-manager-and-pod-analyst:3}
What t-statistic does a three-year daily record with Sharpe ratio 2 have?
\end{exercise}

\begin{exercise}[$\star\star$]\label{exo:in:portfolio-manager-and-pod-analyst:4}
Why is the deal worth \$3.18 million in expectation to a manager with no skill, and who pays for it?
\end{exercise}

\begin{exercise}[$\star\star$]\label{exo:in:portfolio-manager-and-pod-analyst:5}
Why does the certainty equivalent cross the alternative at a higher Sharpe ratio than the expected pay?
\end{exercise}

\begin{exercise}[$\star\star$]\label{exo:in:portfolio-manager-and-pod-analyst:6}
An analyst is offered a sub-portfolio of \$50 million with a 10\% share of its profit, or a \$300\,000 salary with a
discretionary bonus. What must she know to choose?
\end{exercise}

\begin{exercise}[$\star\star\star$]\label{exo:in:portfolio-manager-and-pod-analyst:7}
\emph{Coding.} Rerun the example with the stop moved from 7.5\% to 10\%. What happens to the expected pay and the stop
probability at $S=0$ and $S=1$?
\end{exercise}

\begin{exercise}[$\star\star\star$]\label{exo:in:portfolio-manager-and-pod-analyst:8}
\emph{Find the flaw.} ``She ran a pod for three years with a Sharpe ratio of 1.2 and was never stopped: she is clearly
skilled, and we should give her twice the capital.''
\end{exercise}

\section{Problem: The Deal}

\begin{problem}[{Weekend problem --- the deal}]\label{pb:in:portfolio-manager-and-pod-analyst:1}
A senior analyst at an asset manager, paid \$1.5 million a year, is offered a pod at a platform on the example's terms.
She believes her true Sharpe ratio is between 0.5 and 1.

\medskip\noindent\textbf{Part I --- The roles and the deal.}
\begin{enumerate}
\item Define the \omterm{def:in:portfolio-manager-and-pod-analyst:roles}{portfolio manager}, the \omterm{def:in:portfolio-manager-and-pod-analyst:roles}{sub-portfolio manager} and the \omterm{def:in:portfolio-manager-and-pod-analyst:roles}{pod analyst}.
\item Describe the payout structure the platform's brochure states.
\item What does the platform do with capital, and why is that its main tool?
\item What is the de-risking ladder, and which book defines it?
\item What costs are charged before the payout?
\end{enumerate}

\medskip\noindent\textbf{Part II --- The model.}
\begin{enumerate}[resume]
\item State the deal's payoff with carried losses and a salary advance.
\item Why is it a call, and why a knock-out?
\item State the example's terms.
\item Give the expected pay and stop probability at $S=0$, 1 and 2.
\item What is the certainty equivalent, and at which risk aversion and wealth is it computed?
\end{enumerate}

\medskip\noindent\textbf{Part III --- Records and alternatives.}
\begin{enumerate}[resume]
\item How long must a Sharpe ratio of 1 run to be distinguishable from zero?
\item What does the hiring platform read besides the number?
\item How does a \omterm{def:in:portfolio-manager-and-pod-analyst:roles}{portfolio manager}'s deal at an asset manager differ?
\item What do the survey's wages show, and what do they miss?
\item How does an analyst move toward running capital?
\end{enumerate}

\medskip\noindent\textbf{Part IV --- The verdict.}
\begin{enumerate}[resume]
\item State the \emph{named result}: the true Sharpe ratio at which the expected three-year pay exceeds the salaried
alternative, and the stop probability there.
\item What happens to the answer with her risk aversion?
\item Over her range of beliefs, what are the expected pay and stop probability?
\item What should she negotiate?
\item In two sentences, advise her.
\end{enumerate}
\end{problem}

\section{Interview questions}

\begin{interviewq}[$\star$ \iqroles{trader, researcher}]\label{iq:in:portfolio-manager-and-pod-analyst:1}
Your pod is down 4\% from its peak with a stop at 7.5\%. What do you do?
\end{interviewq}

\begin{interviewq}[$\star$ \iqroles{researcher}]\label{iq:in:portfolio-manager-and-pod-analyst:2}
How would you tell a \omterm{def:in:portfolio-manager-and-pod-analyst:roles}{portfolio manager} that your best idea no longer works?
\end{interviewq}

\begin{interviewq}[$\star\star$ \iqroles{trader}]\label{iq:in:portfolio-manager-and-pod-analyst:3}
Why might a manager with carried losses take more risk, and how would a platform stop it?
\end{interviewq}

\begin{interviewq}[$\star\star$ \iqroles{researcher, trader}]\label{iq:in:portfolio-manager-and-pod-analyst:4}
How much of your pod's return is market beta, and how would you measure it?
\end{interviewq}

\begin{interviewq}[$\star\star$ \iqroles{trader}]\label{iq:in:portfolio-manager-and-pod-analyst:5}
What would you want to know about a candidate manager's three-year record before hiring them?
\end{interviewq}

\begin{interviewq}[$\star\star\star$ \iqroles{researcher, trader}]\label{iq:in:portfolio-manager-and-pod-analyst:6}
Value the manager's option on one year: the pod's profit is normal with mean zero and standard deviation \$30 million, the
payout 15\%. What is the expected payout?
\end{interviewq}
